\documentclass[a4paper]{article}

\usepackage[english]{babel}
%\usepackage[utf8x]{inputenc}
\usepackage{amsmath}
\usepackage{graphicx}
\usepackage[colorinlistoftodos]{todonotes}
\usepackage[all]{xy}
\usepackage{dsfont}
\usepackage{amssymb}
\usepackage{bbold}
\usepackage{amsthm}

\usepackage{tikz}
\usepackage{verbatim}
\usetikzlibrary{matrix}
\usetikzlibrary{cd}

\usepackage[autostyle=true]{csquotes}
\usepackage{pdfpages}
\usepackage{hyperref}
\usepackage{fancyhdr}



\newtheorem{theorem}{Theorem}[section]
\newtheorem{proposition}{Proposition}[theorem]
\newtheorem{lemma}{Lemma}[theorem]
\newtheorem{corollary}{Corollary}[theorem]
\newtheorem*{definition}{Definition}


\newcommand{\norm}[1]{\left\lVert#1\right\rVert}

\newcommand{\Rn}{\mathbb{R}^n}


\newcommand{\too}[1]{\overset{#1}{\longrightarrow}}
\newcommand{\hookar}[1]{\overset{#1}{\hookrightarrow}}
\newcommand{\colim}[1]{ \underset{#1}{\texttt{colim}} \text{ } }
\newcommand{\fset}[1]{ \{ #1_0 , \dots , #1_n \} }
\newcommand{\ddt}{\frac{d}{dt}}


\newcommand{\N}{\mathbb{N}}
\newcommand{\Z}{\mathbb{Z}}







\fancyhead[L]{\includegraphics[scale=0.1]{UFT.png}}

\begin{document}


\thispagestyle{fancy}
\vspace*{1cm}

\begin{minipage}[r]{12cm}
    \begin{flushright}
\begin{tikzpicture}
\draw[green, very thick] (0,0) rectangle (3,2);
\end{tikzpicture}
    \end{flushright}
\end{minipage}

\begin{center}
\Large{\textbf{Control 1\\ Álgebra Lineal\\ }}
\end{center}
\vspace{1cm}
\begin{minipage}[r]{15cm}
\begin{flushleft}

Fecha: 15 de abril 2024 \\%
Duración: 120 minutos  \\
Profesores: Camila Muñoz.\\\hspace{1.7cm} Adrián Silva.\\
\end{flushleft}
\end{minipage}
\vspace{2cm}
\begin{center}
\Large{\textbf{Facultad de Ingeniería}}
\end{center}
\vspace{1cm}
\textbf{En caso de que un alumno sea sorprendido en actitudes deshonestas, ya sea intentando o
dejándose copiar, será calificado con nota mínima 1.0 y perderá toda posibilidad de eliminar
dicha nota. Además, esta conducta quedará archivada como antecedente en su carpeta y
podrá ser causal de eliminación.
Así mismo, está prohibido salir de la sala, usar o mantener sobre la mesa teléfonos celulares
o cualquier tipo de aparato electrónico, a no ser que el profesor lo requiera.}\\



\vspace{5cm}
\makebox[12cm]{\textbf{Nombre Alumno}:\ \hrulefill} 
\newpage
\thispagestyle{empty}
\section*{Preguntas}
\begin{enumerate}


\item {\bf (12 puntos)} Una matriz $U\in M_n(\mathbb{R})$ se dice ortogonal si $U^tU=I_n$.
\begin{enumerate}
    \item Sea $\theta\in \mathbb{R}$ y $G(\theta)\in M_2(\mathbb{R})$ definida por $G(\theta)=\begin{pmatrix} \cos(\theta) & -\sin(\theta)\\ \sin(\theta) & \cos(\theta) \end{pmatrix}$. Pruebe que $U$ es ortogonal para todo $\theta$.
    \item Si $U\in M_n(\mathbb{R})$ es ortogonal, pruebe usando propiedades de determinantes que $U$ es invertible y luego muestre que su inversa tambi\'en es or\-to\-go\-nal.
    \item Si $H, P\in M_n(\mathbb{R})$ son matrices ortogonales, pruebe que el producto $HP$ tambi\'en lo es.
\end{enumerate}
\newpage

\newpage
\item {\bf (12 puntos)} Para las matrices $X$ e $Y$ resuelva el sistema
\begin{align*}
    A^4X-2Y^t&=AB^5\\
    X^t+Y(B^t)^5&=A^t
\end{align*}
donde $A=\begin{pmatrix}0 & 1 \\ -1 & 0 \end{pmatrix}$, y $B=\begin{pmatrix} 3 & 2 \\ -3 & -2 \end{pmatrix}$.
    
\newpage
\newpage
    
%\newpage

%\newpage
\item {\bf (12 puntos)} Utilizando propiedades de determinantes, determine todos los valores de $x$, si se sabe que 

\[\begin{vmatrix} 1 & x  & 3 & 0 \\ 1 & -1 & 0& 1\\ 0 & 0 & -1& -2\\ -x & -1 & 1 &  0  \end{vmatrix}=-9. \]
\newpage
    \item {\bf (12 puntos)} Calcule la inversa de
\[\begin{pmatrix} 1 & 1  & 2  \\ 1 & -1 & 2\\ 0 & 2 & -1 \end{pmatrix}, \]
por el m\'etodo $A^{-1}=\dfrac{1}{\det(A)}\operatorname{adj}(A)$.
\newpage
    \item {\bf (12 puntos)} Usando el teorema del rango, determine los valores de $\alpha$ de modo que el sistema
    \begin{align*}
\alpha x + 3y &= -\alpha \\
\alpha x + 2y +z &= 0\\
2\alpha x + 5y +(\alpha -1) z &= 1-\alpha 
\end{align*}

    tenga solución única, infinitas soluciones o ninguna solución. Para $\alpha=3$, determine el conjunto soluci\'on del sistema.
        
        









\end{enumerate}


\end{document}